The Equational Theories Project recently determined all $22{,}033{,}636$
implications among $4{,}694$ equational laws on magmas, of which $8{,}178{,}279$
(37.1\%) are true.  We investigate whether this implication structure can be
distilled into a small set of interpretable decision rules.  We introduce a
structural classification of equational laws into five \emph{forms} --- trivial,
singleton, absorbing, standard, and general --- and prove that form-based rules
alone decide 39.8\% of all non-reflexive implications with zero error.
Specifically, we show that every absorbing law implies every other law, and
that no general-form law implies any standard-form or absorbing law.  We
further establish a \emph{signature direction} principle: if the left-hand side
of a law $E_1$ contains operations while that of $E_2$ does not, then $E_1$
does not imply~$E_2$, covering an additional 22\% of pairs with near-perfect
accuracy.  For the remaining cases, we prove that the difference in distinct
variable counts is the strongest single predictor of implication.  Composing
these rules into an ordered decision procedure and supplementing with
counterexample magma verification, we achieve 88.9\% accuracy on random samples
and 90\% on a benchmark test set, substantially exceeding the 62.9\%
always-false baseline.  Our results demonstrate that the equational implication
graph, despite its computational origins, admits a surprisingly compact
description through structural algebraic invariants.
